\section{Conclusion}
\label{sec:conclusion}

Complex annotation can be bootstrapped from a guideline and a few examples: without fine-tuning, an LLM that is given the guideline as a chain of steps, tools over the project's resources, and a loop that revises them against a handful of gold sentences, annotates UMR better than the fine-tuned systems of its shared task on the same documents. A long guideline alone is not enough, and neither is a one-time compression of it: the gains come from the harness, its decomposition into steps and its tools, and from recursive self-improvement with a fine-grained judge, which plays the role that moderation plays in a human annotation project. In that loop the builder model learned to edit the guideline, to write the tools the guideline presupposes, and to turn recurring decisions into code.

The same loop shows its limits. What it learns are the conventions of the gold it sees: it learned which TimeML events to link, and, from a sparser corpus, to link less; senses and arguments improved where a tool supplied the frames, but judgments that no resource supports, aspect, coreference and temporal relations, moved little; and its own suggestions are dominated by inconsistencies of the gold. This points to a different division of labor in annotation projects. A guideline can be an executable object that is tested and revised against data from the first annotated sentences on; the model maintains the instructions, tools and checks, and people decide the conventions and adjudicate the inconsistencies it surfaces. Bootstrapping a new scheme or language then starts from a working annotator rather than from an empty corpus.
